% ============================================================
%  PCL-02 Implementation Paper
%  Title: Measuring and Mitigating Context-Window Degradation
%         in Long-Document Research Synthesis: 
%         A Controlled Retrieval Experiment
%  Author: Mayank Waila
%  Institute: Indian Institute of Information Technology Allahabad
% ============================================================

\documentclass[sigconf,nonacm]{acmart}

% ── packages ─────────────────────────────────────────────────────────────────
\usepackage{booktabs}
\usepackage{graphicx}
\usepackage{listings}
\usepackage{xcolor}
\usepackage{microtype}
\usepackage{url}
\usepackage{amsmath}
\usepackage{multirow}
\usepackage{subcaption}

% ── code style ───────────────────────────────────────────────────────────────
\lstset{
  basicstyle=\ttfamily\footnotesize,
  breaklines=true,
  frame=single,
  numbers=left,
  numberstyle=\tiny,
  keywordstyle=\color{blue},
  commentstyle=\color{gray},
  stringstyle=\color{red},
  language=Python
}

% ── ACM metadata ─────────────────────────────────────────────────────────────
\acmJournal{PACMPL}

\begin{document}

\title{Measuring and Mitigating Context-Window Degradation in\\
Long-Document Research Synthesis:\\
A Controlled Retrieval Experiment}

\author{Mayank Waila}
\affiliation{%
  \institution{Indian Institute of Information Technology Allahabad}
  \city{Prayagraj}
  \country{India}}
\email{mayankwaila@iiita.ac.in}

\begin{abstract}
Large language models (LLMs) are increasingly deployed for research synthesis
over long documents, yet their performance degrades when relevant evidence is
positioned away from the context boundaries—the ``lost-in-the-middle'' phenomenon
documented by Liu \emph{et al.}~\cite{liu2023lost}.
This paper presents a controlled, reproducible experimental investigation of
context-window degradation in a long-document retrieval task.
We implement a U-shaped attention-weight simulation calibrated to match
empirically observed LLM behaviour, then compare four retrieval strategies:
(1)~full-context single-pass retrieval (Baseline~1), (2)~non-overlapping chunked
sequential retrieval (Baseline~2), (3)~overlapping chunked retrieval (Baseline~3),
and (4)~a proposed two-pass hierarchical retrieval method.
Over 200 independent trials per condition across five evidence-position levels
(0\%, 25\%, 50\%, 75\%, 100\% into a 20-chunk document), we find that the full-context
baseline exhibits a \textbf{10 percentage-point recall drop} at non-boundary positions
(Recall@1 = 0.900--0.905 vs.\ 1.000 at the boundaries), while all three
chunking-based methods recover full recall at every position.
An ablation over attention-decay strength confirms that degradation is robust
above a strength threshold of 1.0, and plateaus thereafter.
Our code, configuration, and raw results are released in a public repository,
enabling full reproduction of all reported numbers.
\end{abstract}

\begin{CCSXML}
<ccs2012>
  <concept>
    <concept_id>10010147.10010178.10010187</concept_id>
    <concept_desc>Computing methodologies~Natural language processing</concept_desc>
    <concept_significance>500</concept_significance>
  </concept>
  <concept>
    <concept_id>10002951.10003317.10003338</concept_id>
    <concept_desc>Information systems~Information retrieval</concept_desc>
    <concept_significance>300</concept_significance>
  </concept>
</ccs2012>
\end{CCSXML}

\ccsdesc[500]{Computing methodologies~Natural language processing}
\ccsdesc[300]{Information systems~Information retrieval}

\keywords{context-window degradation, long-context LLMs,
  retrieval-augmented generation, hierarchical retrieval,
  lost-in-the-middle, research synthesis}

\maketitle

% ─────────────────────────────────────────────────────────────────────────────
\section{Introduction}
\label{sec:intro}
% ─────────────────────────────────────────────────────────────────────────────

The capacity of large language models (LLMs) to process tens of thousands of
tokens has opened compelling applications in academic research synthesis, where
a system must traverse long survey articles, multi-document corpora, or book-length
texts to locate, compare, and integrate evidence~\cite{surveylong2024}.
A naïve view holds that wider context windows eliminate the retrieval problem
entirely: if the model can read the whole document, it should find any piece of
evidence within it.
Empirical work has repeatedly refuted this view.
Liu \emph{et al.}~\cite{liu2023lost} showed that GPT-3.5 and Claude~1.3 exhibit
substantially lower accuracy when the relevant passage is placed in the
\emph{middle} of a long input, producing the now well-known \emph{lost-in-the-middle}
phenomenon.
Bai \emph{et al.}~\cite{bai2024longbench} and Hsieh \emph{et al.}~\cite{hsieh2024ruler}
independently confirmed and extended these findings, demonstrating that
long-context performance can degrade even when the full text fits within a
model's nominal token budget.

My previous literature review (Paper~1) synthesised this body of work and
identified context-window degradation as a central reliability problem for
long-document research synthesis, noting that the phenomenon had been characterised
but not yet systematically addressed under a controlled, apples-to-apples
retrieval protocol with open, reproducible code.
This paper is the direct continuation of that review.
It narrows the scope to a single tractable research question (Section~\ref{sec:rq}),
implements four retrieval methods in runnable Python code (Section~\ref{sec:methods}),
and reports all numbers from actual experimental runs (Section~\ref{sec:results}).

\subsection{Research Question and Hypothesis}
\label{sec:rq}

\begin{quote}
\textbf{RQ:} \emph{Does the positional placement of target evidence within a
long document significantly degrade retrieval accuracy under a simulated LLM
attention model, and can chunked or hierarchical retrieval strategies recover
this loss?}
\end{quote}

\noindent\textbf{H1 (Degradation):} Retrieval recall in a full-context single-pass
strategy will be significantly lower when evidence is placed in the middle of
the document compared to the boundaries.

\noindent\textbf{H2 (Mitigation):} Chunked and hierarchical retrieval strategies
that subdivide the context will maintain consistently high recall regardless of
evidence position, because the effective window over which attention must span
is reduced in each retrieval call.

\subsection{Contributions}

\begin{enumerate}
  \item A reproducible simulation framework that operationalises the
        lost-in-the-middle phenomenon via a parameterised attention-decay curve.
  \item A controlled comparison of four retrieval strategies across five
        evidence positions, with 200 independent trials per condition.
  \item An ablation study over attention-decay strength, isolating the
        threshold at which degradation becomes harmful.
  \item Full open-source release: code, configuration, and raw CSV results.
\end{enumerate}

% ─────────────────────────────────────────────────────────────────────────────
\section{Background and Related Work}
\label{sec:background}
% ─────────────────────────────────────────────────────────────────────────────

\subsection{The Lost-in-the-Middle Phenomenon}

Liu \emph{et al.}~\cite{liu2023lost} evaluated GPT-3.5-Turbo and Claude~1.3 on
multi-document question answering while systematically varying the position of
the gold document within the input.
They found a U-shaped performance curve: accuracy was highest when the relevant
document appeared first or last, and dropped by up to 20 percentage points when
it appeared in the middle of a 20-document context.
This positional sensitivity was robust across both models and both tasks
(QA and key-value retrieval).

\subsection{Long-Context Benchmarks}

LongBench~\cite{bai2024longbench} and RULER~\cite{hsieh2024ruler} provide
standardised evaluation suites that reveal similar patterns at the benchmark
scale.
LongBench includes six categories of long-context tasks including single-document
QA, multi-document QA, and summarisation.
RULER probes recall of needle-in-a-haystack facts, multi-key retrieval, and
chain-of-reasoning tasks up to 128k tokens.
Both suites show that models with higher nominal context limits do not
proportionally outperform smaller-window models on retrieval-intensive sub-tasks.

\subsection{Retrieval-Augmented Generation (RAG)}

RAG~\cite{lewis2020rag} was originally proposed to inject retrieved document
chunks into a generation context, sidestepping the long-context problem entirely.
Chunked retrieval creates a short, focused prompt for each generation call,
effectively resetting the attention window.
Hierarchical and coarse-to-fine variants have been explored in the literature
but rarely evaluated with controlled positional benchmarks~\cite{gao2023retrieval}.

\subsection{Context Compression}

LLMLingua~\cite{jiang2023llmlingua} and similar approaches compress the prompt
prior to model ingestion, retaining a higher information density per token.
While promising, compression methods risk dropping critical evidence,
particularly when that evidence uses rare terminology—precisely the signal
tokens we use in our synthetic experiment.

\subsection{Position of This Work}

The body of work above identifies and characterises the problem but leaves open
a systematic, controlled comparison of simple chunking strategies against the
full-context baseline under a consistent evaluation protocol.
This paper fills that gap.

% ─────────────────────────────────────────────────────────────────────────────
\section{Experimental Setup}
\label{sec:methods}
% ─────────────────────────────────────────────────────────────────────────────

\subsection{Synthetic Document Generation}
\label{sec:data}

We construct synthetic long documents to allow precise control over evidence
placement—a capability unavailable in natural corpora where document structure
confounds positional effects.

Each document consists of $C = 20$ text chunks of 60 words each.
One chunk (the \emph{evidence chunk}) contains three unique signal tokens
(\texttt{SIGNAL\_X}, \texttt{SIGNAL\_Y}, \texttt{SIGNAL\_Z}) not present in
the distractor vocabulary.
All remaining chunks contain random tokens drawn from a 800-word distractor vocabulary.
The query consists of two of the three signal tokens.

Evidence placement is controlled by the parameter $p \in \{0, 5, 10, 15, 19\}$,
corresponding to positions 0\%, 25\%, 50\%, 75\%, and 100\% into the document
(i.e., Start, Early, Middle, Late, End).
For each position and method, $T = 200$ independent trials are run with distinct
random seeds, yielding a total of $4 \times 5 \times 200 = 4{,}000$ retrievals.
All random seeds are stored in \path{config_v2.json} and reproduced exactly
from the seed formula: $\text{seed}(p_i, t) = 42 + p_i \times 10{,}000 + t$.

\subsection{Attention-Decay Simulation}
\label{sec:attn}

Since we cannot run a full LLM in a compute-limited setting without API access,
we simulate the primacy-recency attention curve empirically observed by
Liu \emph{et al.}~\cite{liu2023lost}.
Following their reported ~20 pp drop at the middle, we parameterise:

\begin{equation}
  w(p, N, \lambda) = \max\!\left(0.05,\; 1 - \lambda \cdot \sin^2\!\!\left(\frac{\pi p}{N-1}\right)\right)
  \label{eq:attn}
\end{equation}

\noindent where $p$ is the chunk index (0-indexed), $N$ is the total number of
chunks in scope, and $\lambda$ is the \emph{decay strength} hyper-parameter.
At $p = 0$ and $p = N-1$ (boundaries), $w = 1.0$.
At $p = \lfloor(N-1)/2\rfloor$ (middle), $w = 1 - \lambda$.
We set $\lambda = 1.8$ for the main experiments, which produces a ~10 pp
performance drop at Early/Middle positions (matching the lower end of Liu \emph{et al.}'s
reported range, appropriate for our bag-of-words retriever baseline rather than
a full LLM).

The final score for chunk $i$ is:
\begin{equation}
  s_i = \text{BoW}(c_i, q) \cdot w(i, N, \lambda) + \varepsilon, \quad \varepsilon \sim \mathcal{N}(0, 0.015^2)
  \label{eq:score}
\end{equation}

\noindent where $\text{BoW}(c, q)$ is the token-overlap recall score
(query tokens found in chunk $c$, divided by total query tokens), and $\varepsilon$
is Gaussian noise simulating stochastic LLM behaviour.

\subsection{Retrieval Methods}
\label{sec:methods-detail}

We compare four retrieval strategies:

\paragraph{Baseline 1 — Full-Context (FC-B1).}
All $C = 20$ chunks are scored simultaneously using Equation~\ref{eq:score} with the
global window size $N = C$.
This is the ``single-pass long context'' baseline corresponding to naïve full-document
ingestion.
\emph{Expected weakness:} chunks in the middle receive low $w(\cdot)$ weights,
causing the model to miss central evidence.

\paragraph{Baseline 2 — Chunked Sequential (CS-B2).}
The document is processed in non-overlapping windows of size $W = 5$.
Within each window, attention decay is applied \emph{locally} (positions 0–4
within the window) with strength $\lambda' = 0.8$.
The highest-scoring chunk across all windows is returned.
\emph{Hypothesis:} smaller windows reduce degradation at the cost of potential
boundary artefacts.

\paragraph{Baseline 3 — Overlapping Chunks (OC-B3).}
Like CS-B2 but with a stride of 2 (windows overlap by half), ensuring that no
chunk appears only in the middle of its window.
This directly addresses the ``chunk-boundary problem'' where evidence straddles
two non-overlapping windows.

\paragraph{Proposed — Hierarchical Two-Pass (H2P).}
The proposed method operates in two stages:

\begin{enumerate}
  \item \textbf{Coarse pass:} Each non-overlapping window is represented by the
    maximum BoW score of its first 12-token chunk summaries.
    The top 2 windows are selected.
  \item \textbf{Fine pass:} Chunks within the selected windows are re-scored
    with a much lower decay strength ($\lambda_\text{fine} = 0.25$), simulating
    a short, focused context where positional effects are minimal.
\end{enumerate}

The two-pass design is motivated by the intuition that a coarse filter can
identify \emph{candidate regions} without full attention, after which a short,
focused re-read applies nearly uniform attention.

\subsection{Metrics}

We report:
\begin{itemize}
  \item \textbf{Recall@1}: proportion of trials in which the top-ranked chunk is
    the evidence chunk.
  \item \textbf{Mean F1}: average token-overlap F1 between the retrieved chunk
    and the evidence chunk (1.0 for a correct retrieval, 0.0 otherwise, since
    signal tokens are unique).
\end{itemize}

\subsection{Reproducibility}

All experiments are implemented in a single Python file
(\path{experiment_v2.py}) with no external dependencies beyond the Python
standard library. The file is 170 lines of commented code.
Configuration (seeds, chunk counts, positions, trials) is serialised to
\path{config_v2.json}. Outputs are written to \path{results_v2.csv}
and~\path{ablation_strength.csv}. Running \path{python3 experiment_v2.py}
from the repository root reproduces all numbers in this paper exactly.

% ─────────────────────────────────────────────────────────────────────────────
\section{Results}
\label{sec:results}
% ─────────────────────────────────────────────────────────────────────────────

\subsection{Main Results: Recall@1 by Position and Method}

Table~\ref{tab:main} reports Recall@1 for all four methods across five evidence
positions, based on $T = 200$ trials per cell.

\begin{table}[t]
  \centering
  \caption{Recall@1 (200 trials per cell, seed 42). Bold = best per row.
           FC-B1: Full-Context; CS-B2: Chunked Sequential;
           OC-B3: Overlapping Chunks; H2P: Hierarchical Two-Pass.}
  \label{tab:main}
  \begin{tabular}{lcccc}
    \toprule
    Evidence Position & FC-B1 & CS-B2 & OC-B3 & H2P \\
    \midrule
    Start (0\%)       & \textbf{1.000} & \textbf{1.000} & \textbf{1.000} & \textbf{1.000} \\
    Early (25\%)      & 0.900 & \textbf{1.000} & \textbf{1.000} & \textbf{1.000} \\
    Middle (50\%)     & 0.905 & \textbf{1.000} & \textbf{1.000} & \textbf{1.000} \\
    Late  (75\%)      & \textbf{1.000} & \textbf{1.000} & \textbf{1.000} & \textbf{1.000} \\
    End   (100\%)     & \textbf{1.000} & \textbf{1.000} & \textbf{1.000} & \textbf{1.000} \\
    \midrule
    \textbf{Avg.}     & 0.961 & \textbf{1.000} & \textbf{1.000} & \textbf{1.000} \\
    \bottomrule
  \end{tabular}
\end{table}

\paragraph{Findings.}
H1 is confirmed: the full-context baseline (FC-B1) exhibits a statistically
meaningful recall drop of approximately 10 percentage points at the Early (25\%)
and Middle (50\%) positions (Recall@1 = 0.900 and 0.905, respectively), compared
to perfect recall at the boundaries.
This replicates the direction of the effect documented by Liu \emph{et al.}~\cite{liu2023lost}
in controlled conditions.

H2 is also confirmed: all three chunking-based strategies (CS-B2, OC-B3, H2P)
achieve Recall@1 = 1.000 at every position.
Subdividing the retrieval context into smaller windows eliminates the positional
degradation observed in the full-context condition.

\paragraph{Why does the full-context baseline still score 0.90 rather than lower?}
Under our simulation, the BoW signal is strong (3 unique tokens not in the
distractor vocabulary), so the base score is 1.0 for the evidence chunk and 0.0
for all distractors.
The attention-decay weight at the middle position is approximately $1 - 1.8 = -0.8$,
clipped to 0.05.
Gaussian noise $\mathcal{N}(0, 0.015^2)$ then occasionally causes a distractor
to outscore the evidence chunk—accounting for the ~10\% error rate.
In a real LLM where distractors also contribute partial token overlaps, the gap
would be expected to be larger, consistent with the 20 pp drops reported by
Liu \emph{et al.}\ for GPT-3.5.

\subsection{Ablation: Attention-Decay Strength}

Table~\ref{tab:ablation} shows how Recall@1 at the Middle (50\%) position varies
as a function of $\lambda$ in Equation~\ref{eq:attn}.

\begin{table}[t]
  \centering
  \caption{Ablation: Recall@1 of FC-B1 at Middle (50\%) vs.\ attention-decay
           strength $\lambda$ (200 trials, seed 99999).}
  \label{tab:ablation}
  \begin{tabular}{cc}
    \toprule
    Decay strength $\lambda$ & Recall@1 \\
    \midrule
    0.0 (uniform attention) & 1.000 \\
    0.5                     & 1.000 \\
    1.0                     & 0.920 \\
    1.5                     & 0.920 \\
    1.8 (main experiment)   & 0.920 \\
    2.0                     & 0.920 \\
    2.5                     & 0.920 \\
    3.0                     & 0.920 \\
    \bottomrule
  \end{tabular}
\end{table}

\paragraph{Findings.}
Degradation is essentially absent below $\lambda = 1.0$, then plateaus at
Recall@1 = 0.920 for all $\lambda \geq 1.0$.
This plateau reflects the signal strength in our synthetic setup: once the
attention weight at the middle drops below the signal-to-noise threshold, further
increases in $\lambda$ cannot worsen performance because noise alone accounts for
the remaining errors.
Importantly, this plateau means that the magnitude of our main result
(Recall@1 $\approx$ 0.90--0.92 for FC-B1 at non-boundary positions) is
\emph{robust} to the exact value of $\lambda \geq 1.0$ chosen to calibrate the
simulation.

\subsection{Variance and Statistical Reliability}

With $T = 200$ binary trials per cell, the 95\% confidence interval for a
proportion $\hat{p}$ is approximately $\hat{p} \pm 1.96\sqrt{\hat{p}(1-\hat{p})/T}$.
For $\hat{p} = 0.90$: CI $= [0.858, 0.942]$.
For $\hat{p} = 1.00$: CI $= [1.000, 1.000]$ (degenerate; no failures were
observed in 200 trials for any chunking method at any position).
The gap between FC-B1 and the chunking methods at the Early/Middle positions
is therefore statistically significant at the 95\% level.

% ─────────────────────────────────────────────────────────────────────────────
\section{Analysis and Discussion}
\label{sec:discussion}
% ─────────────────────────────────────────────────────────────────────────────

\subsection{Why Chunking Recovers Recall}

The fundamental reason all three chunking strategies outperform the full-context
baseline is that they reduce the \emph{effective} $N$ over which
Equation~\ref{eq:attn} is applied.
In the chunked sequential method, $N = W = 5$ within each window; the evidence
chunk can appear at most at position 4 within its window (the ``end'' of the
window), and the worst-case weight is $1 - 0.8 \cdot \sin^2(\pi \cdot 4/4) = 1.0$
(also a boundary).
In the hierarchical method, the fine-pass operates on 5-chunk windows with
$\lambda_\text{fine} = 0.25$, giving a worst-case middle weight of $1 - 0.25 = 0.75$—
still high enough to dominate the noise.

This confirms the theoretical account proposed by Liu \emph{et al.}~\cite{liu2023lost}
and echoed in RAG surveys~\cite{gao2023retrieval}: short contexts are simply
easier for the retrieval mechanism to traverse, regardless of the retriever's
implementation details.

\subsection{The Chunk-Boundary Problem}

The initial experiment (v1, using non-overlapping windows of size 5) revealed
an interesting edge case: for window sizes 3 and 7, recall at the middle dropped
to 0.08 (from 1.00 for sizes 2, 5, and 10).
Post-hoc analysis confirmed a \emph{chunk-boundary effect}: when the evidence
chunk sits at the boundary between two non-overlapping windows, the first window
truncates the evidence and the second window's context starts \emph{after} the
key signal tokens.
The overlapping-chunks method (OC-B3) with stride 2 fully addresses this:
every chunk appears in at least two windows, ensuring that the signal tokens
appear near the \emph{start} of at least one window.
This is a nuanced failure mode not discussed by Liu \emph{et al.} and
represents a novel finding of this study.

\subsection{Limitations}

\paragraph{Simulation vs.\ real LLM.}
Our attention-decay model is a first-order approximation calibrated to published
results, not a measurement of any specific model.
Real LLMs are autoregressive; their attention patterns depend on the entire
sequence of previously processed tokens, not just position.
Future work should validate these findings with API-accessible LLMs using
controlled needle-in-a-haystack prompts.

\paragraph{Synthetic data.}
Our signal tokens are unique and high-contrast.
In real research synthesis, evidence may be semantically similar to surrounding
text (e.g., a cited figure surrounded by discussion of the same topic), making
retrieval harder for all methods.
Repeating this experiment with dense retrievers (bi-encoders) on real academic
datasets (NarrativeQA, QuALITY, SCROLLS) would strengthen external validity.

\paragraph{No latency or cost modelling.}
Chunked methods make multiple retrieval calls, which incurs higher latency and
API cost.
A practical deployment trade-off analysis is left for future work.

\subsection{Implications for Research Synthesis Systems}

The results have direct practical implications.
A research assistant system that ingests a full 50-page survey article in a
single LLM call and asks ``What did Liu et al.\ find about middle-passage
retrieval?'' may reliably fail if the cited study appears on pages 20--30.
Replacing the single-pass ingestion with a hierarchical retrieval layer—coarse
section-level filtering followed by fine paragraph-level re-ranking—would
recover this recall at a modest computational cost.
The H2P method demonstrated here is an existence proof that such recovery is
achievable even with simple BoW scoring.

% ─────────────────────────────────────────────────────────────────────────────
\section{Human--AI Collaboration Log}
\label{sec:ai}
% ─────────────────────────────────────────────────────────────────────────────

This section documents the AI tools used, following the transparency requirements
of the course assignment.

\begin{enumerate}
  \item \textbf{Research question formulation.}
    The research question (Section~\ref{sec:rq}) was derived from my own reading
    of the five primary sources cited in Paper~1.
    AI was used to check whether the question was sufficiently narrow and
    experimentally testable (Prompt: ``Is this question tractable with a
    compute-budget-free Python simulation?'').
    AI confirmed feasibility; the question itself was not generated by AI.

  \item \textbf{Attention-decay equation.}
    I specified that I wanted a U-shaped curve with boundary weights of 1.0 and
    a user-controllable depth.
    AI suggested the $\sin^2$ parameterisation.
    I independently verified that it matches the shape reported in Liu et al.\
    Figure~2 (U-shaped with ~20 pp drop) by computing values at five positions
    and cross-checking with the reported numbers.
    The final equation was written by me; the shape suggestion came from AI.

  \item \textbf{Code scaffolding.}
    Initial Python scaffold was generated with AI assistance (vibe coding).
    I verified, debugged, and rewrote approximately 60\% of the code, including:
    the trial-seeding formula, the overlapping-window stride logic, the two-pass
    hierarchical method, and the CSV output format.
    All bugs introduced by AI (e.g., off-by-one error in chunk-boundary indexing)
    were caught by running the code and inspecting the v1 ablation anomaly.

  \item \textbf{Paper writing.}
    Section drafts were written by me.
    AI was used to suggest sentence-level rephrasing in two paragraphs of the
    Introduction and to verify that the LaTeX table syntax compiled correctly.
    All claims, numbers, and analysis are my own.

  \item \textbf{Reference management.}
    BibTeX entries were checked against ACM DL and arXiv DOIs manually.
    AI was not used to fabricate references.
\end{enumerate}

\noindent\textbf{Decision ownership statement.}
Every technical decision in this paper—choice of decay function, window size,
stride, number of trials, evidence positions, and interpretation of results—was
made by the author with reasoning documented in the code comments and in this
section.
AI output was treated as a suggestion requiring verification, not as authority.

% ─────────────────────────────────────────────────────────────────────────────
\section{Conclusion}
\label{sec:conclusion}
% ─────────────────────────────────────────────────────────────────────────────

This paper presented a controlled, reproducible experimental investigation of
context-window degradation in long-document research synthesis retrieval.
Using a parameterised U-shaped attention-decay simulation calibrated to the
lost-in-the-middle phenomenon, we confirmed that:
\begin{enumerate}
  \item Full-context single-pass retrieval exhibits significant recall degradation
        (~10 pp) when evidence is placed at non-boundary positions (H1 confirmed).
  \item All three chunking-based strategies—non-overlapping, overlapping, and
        hierarchical two-pass—fully recover this recall at every position (H2 confirmed).
  \item Degradation saturates above an attention-decay strength threshold of 1.0,
        making the main result robust to the exact calibration chosen.
  \item Non-overlapping chunking introduces a chunk-boundary failure mode that
        overlapping and hierarchical methods avoid—a novel finding not previously
        documented in the literature.
\end{enumerate}

The practical takeaway is straightforward: research synthesis systems should
not rely on monolithic full-context ingestion.
Even simple sliding-window chunking—implementable in tens of lines of Python—
eliminates positional retrieval failures.
Hierarchical two-pass retrieval, while equally effective in this controlled
setting, provides an additional degree of freedom (the coarse filter) that may
generalise better to noisy real-world documents where distractors partially
match the query.

\paragraph{Future work.}
Validating these findings with real LLM API calls (GPT-4, Gemini, Claude),
dense bi-encoder retrievers, and natural academic texts (SCROLLS benchmark)
are the most important next steps.
Investigating the interaction between decay strength and document length—
specifically whether the plateau persists at $C \gg 20$—is a secondary priority.

% ─────────────────────────────────────────────────────────────────────────────
\section*{Reproducibility Statement}
% ─────────────────────────────────────────────────────────────────────────────

All results in this paper are reproduced by running:
\begin{lstlisting}[language=bash]
git clone https://github.com/mayankwaila/context-window-degradation
cd context-window-degradation
python3 experiment_v2.py
\end{lstlisting}
No external dependencies beyond the Python 3.8+ standard library are required.
The script writes \path{results_v2.csv}, \path{ablation_strength.csv},
and \path{config_v2.json} to the working directory.
All random seeds, hyper-parameters, and output file paths are listed in
\path{config_v2.json}.

% ─────────────────────────────────────────────────────────────────────────────
\bibliographystyle{ACM-Reference-Format}
\bibliography{references}

\end{document}
